\begin{proof}[Proof of \cref{obj:thm:chebyshev-minimax}]
\begin{enumerate}
\item \emph{The constants.}
By \cref{obj:lem:continuous-chebyshev-endpoint-bound}, applied with the cap \(q_{\max}\), choose constants
\(C_{\mathrm{ep}}(q_{\max})>0\) and \(c_{\mathrm{ep}}(q_{\max})>0\) such that, writing
\(x_q=2/q-1\) and
\[
\lambda(q)=x_q+\sqrt{x_q^2-1},
\]
the following two bounds hold for every \(\beta\ge1\) and every \(0<q\le q_{\max}\): first, every real polynomial \(R\) with \(\deg(R)\le\beta\) and \(\sup_{x\in[-1,1]}|R(x)|\le1\) satisfies
\[
|R(x_q)-R(-1)|\le C_{\mathrm{ep}}(q_{\max})\,\lambda(q)^\beta ;
\]
and second,
\[
T_\beta(x_q)-1\ge c_{\mathrm{ep}}(q_{\max})\,\lambda(q)^\beta .
\]
These constants are chosen before the oversampling ratio is fixed, so set
\[
C_-(q_{\max})=c_{\mathrm{ep}}(q_{\max})^2>0.
\]
Now fix \(c>1\). The Ehlich--Zeller Chebyshev--Lobatto mesh inequality used as the premise of \cref{obj:lem:oversampled-chebyshev-lobatto-norming} is available in the following form: for every pair of nonnegative integers \(\beta,k\), every real polynomial \(R\) with \(\deg(R)\le\beta\) and \(\beta<k\), and every \(x\in[-1,1]\),
\[
 |R(x)|
 \le
 \frac{1}{\cos(\pi\beta/(2k))}
 \max_{0\le j\le k}\left|R\!\left(-\cos\!\left(\frac{\pi j}{k}\right)\right)\right|.
\]
% lean: ehlichZellerMesh
Applying \cref{obj:lem:oversampled-chebyshev-lobatto-norming} with this mesh inequality, choose
\(K(c)>0\) with the property that every real polynomial \(R\) of degree at most \(\beta\) obeys
\(\sup_{x\in[-1,1]}|R(x)|\le K(c)\max_{0\le j\le k}\left|R(-\cos(\pi j/k))\right|\) whenever \(\beta\ge1\) and \(k\ge c\beta\), and set
\[
C_+(c,q_{\max})=\bigl(K(c)C_{\mathrm{ep}}(q_{\max})\bigr)^2>0 .
\]
% lean: chebyshev_minimax

\item \emph{Standing consequences of the hypotheses.}
Fix \(\beta,k,q\) satisfying the hypotheses. The low-budget-regime hypothesis gives \(q>0\), and the low-budget cap of \cref{obj:ass:low-budget-cap} gives
\(q\le q_{\max}<1\); hence \(q\le 1\), \(q<1\), and \(0<q\le 1\). Since \(k\ge c\beta\), \(c>1\), and
\(\beta\ge1\), we have \(\beta\le c\beta\le k\), and in particular \(k\ge1\). Therefore
\cref{obj:lem:chebyshev-schedule-admissible} gives
\[
p^{\mathrm{Ch}}(k,q)\in S_{k,q}.
\]
% lean: chebyshev_schedule_admissible
The exterior parameter is the exponential base in the statement:
\[
\lambda(q)=\frac{(1+\sqrt{1-q})^2}{q}.
\]
Indeed, \(x_q^2-1=(2/q-1)^2-1=4/q^2-4/q=4(1-q)/q^2\), which is nonnegative because \(q<1\), so \(\sqrt{x_q^2-1}=2\sqrt{1-q}/q\); adding \(x_q=(2-q)/q\) gives \(\lambda(q)=\bigl(2-q+2\sqrt{1-q}\bigr)/q\), and expanding \((1+\sqrt{1-q})^2=1+2\sqrt{1-q}+(1-q)=2-q+2\sqrt{1-q}\) identifies the two expressions.
% lean: chebyshev_lambda_eq_rho_div_q
Finally, for any schedule \(p'\in S_{k,q}\) we record the dual quantity of \cref{obj:lem:amplification-dual-norm},
\[
D(p')=\sup\left\{|r(1)-r(0)|:
r\in\mathbb R[x],\ \deg(r)\le\beta,\ \max_{0\le j\le k}|r(p'_j)|\le1\right\},
\]
and note that \(D(p')\ge0\) because the zero polynomial is feasible.

\item \emph{Lower bound at an arbitrary \(p\in S_{k,q}\).}
Since \(\beta\ge1\), \(\beta\le k\), and \(p\in S_{k,q}\), \cref{obj:lem:unbiased-weight-set-nonempty}
makes \(W_\beta(p)\) nonempty; and the strict ordering \(p_0<\cdots<p_k\) required by
\cref{obj:def:budgeted-schedule-class} makes the nodes distinct. Thus
\cref{obj:lem:amplification-dual-norm} applies and gives
\[
A_\beta(p)=D(p)^2 .
\]
% lean: amplification_dual_norm

Consider the test polynomial
\[
r(u)=T_\beta\!\left(\frac{2u}{q}-1\right).
\]
Its degree is at most \(\beta\), being the composition of the degree-\(\beta\) polynomial \(T_\beta\) with an affine map. % lean: chebyshev_affine_natDegree_le
For every rollout node, \(p\in S_{k,q}\) gives \(p_j\ge p_0=0\) and, by monotonicity, \(p_j\le p_k=q\); hence
\[
0\le p_j\le q,
\qquad\text{so}\qquad
\frac{2p_j}{q}-1\in[-1,1],
\]
and since \(|T_\beta|\le1\) on \([-1,1]\) we get \(|r(p_j)|\le1\). % lean: budgetedSchedule_le_endpoint
Therefore \(r\) is feasible in the supremum defining \(D(p)\). Its endpoint values are
\[
r(1)=T_\beta\!\left(\frac2q-1\right)=T_\beta(x_q),
\qquad
r(0)=T_\beta(-1).
\]
The endpoint lower bound from \cref{obj:lem:continuous-chebyshev-endpoint-bound} gives
\[
T_\beta(x_q)-1\ge c_{\mathrm{ep}}(q_{\max})\lambda(q)^\beta .
\]
Since \(x_q>1\) we have \(T_\beta(x_q)\ge1>0\), so \(|T_\beta(x_q)|=T_\beta(x_q)\); and \(|T_\beta(-1)|\le1\). The reverse triangle inequality therefore yields
\[
T_\beta(x_q)-1
\le
|T_\beta(x_q)|-|T_\beta(-1)|
\le |T_\beta(x_q)-T_\beta(-1)|
=|r(1)-r(0)|
\le D(p).
\]
Hence \(D(p)\ge c_{\mathrm{ep}}(q_{\max})\lambda(q)^\beta\ge0\), and squaring both nonnegative sides,
\[
A_\beta(p)=D(p)^2\ge c_{\mathrm{ep}}(q_{\max})^2\lambda(q)^{2\beta}
=
C_-(q_{\max})
\left(\frac{(1+\sqrt{1-q})^2}{q}\right)^{2\beta}.
\]
% lean: chebyshev_amplification_lower

\item \emph{Upper bound at the Chebyshev--Lobatto schedule.}
The schedule \(p^{\mathrm{Ch}}(k,q)\) is admissible by \cref{obj:lem:chebyshev-schedule-admissible}. Together with \(\beta\ge1\) and \(\beta\le k\), \cref{obj:lem:unbiased-weight-set-nonempty} makes \(W_\beta(p^{\mathrm{Ch}}(k,q))\) nonempty; the strict ordering in \(p^{\mathrm{Ch}}(k,q)\in S_{k,q}\) gives distinct nodes. Thus \cref{obj:lem:amplification-dual-norm} gives
\(A_\beta(p^{\mathrm{Ch}}(k,q))=D(p^{\mathrm{Ch}}(k,q))^2\), and it suffices to bound the dual supremum. Let \(r\) be
any polynomial with \(\deg(r)\le\beta\) and
\(\max_{0\le j\le k} |r(p^{\mathrm{Ch}}_j(k,q))|\le1\), and put
\[
R(x)=r\!\left(\frac q2(x+1)\right).
\]
Then \(\deg(R)\le\beta\), again because the inner map is affine. % lean: lobatto_affine_comp_natDegree_le
At the Lobatto abscissae the inner map reproduces the schedule,
\[
\frac q2\left(-\cos\frac{\pi j}{k}+1\right)
=
\frac q2\left(1-\cos\frac{\pi j}{k}\right)
=
p^{\mathrm{Ch}}_j(k,q),
\qquad j=0,\ldots,k,
\]
so \(R(-\cos(\pi j/k))=r(p^{\mathrm{Ch}}_j(k,q))\) and therefore
\[
\max_{0\le j\le k}\left|R\!\left(-\cos\frac{\pi j}{k}\right)\right|\le1 .
\]
The oversampled norming inequality supplied by \cref{obj:lem:oversampled-chebyshev-lobatto-norming} with the displayed Ehlich--Zeller mesh input applies because \(\beta\ge1\), \(k\ge c\beta\), and \(\deg(R)\le\beta\). It gives
\[
|R(x)|\le K(c)\cdot 1=K(c)\qquad (x\in[-1,1]).
\]
Thus \(S=K(c)^{-1}R\) satisfies \(\deg(S)\le\beta\) and \(|S(x)|\le1\) on \([-1,1]\). The inner affine map sends the two relevant points to the two endpoints,
\[
\frac q2\left(x_q+1\right)=\frac q2\cdot\frac2q=1,
\qquad
\frac q2\left(-1+1\right)=0,
\]
so \(R(x_q)=r(1)\) and \(R(-1)=r(0)\), and hence \(S(x_q)=K(c)^{-1}r(1)\) and \(S(-1)=K(c)^{-1}r(0)\). Applying the endpoint upper bound from
\cref{obj:lem:continuous-chebyshev-endpoint-bound} to \(S\) gives
\[
K(c)^{-1}|r(1)-r(0)|
=
|S(x_q)-S(-1)|
\le C_{\mathrm{ep}}(q_{\max})\lambda(q)^\beta,
\]
that is, \(|r(1)-r(0)|\le K(c)C_{\mathrm{ep}}(q_{\max})\lambda(q)^\beta\). Taking the supremum over all such \(r\) gives
\(0\le D(p^{\mathrm{Ch}}(k,q))\le K(c)C_{\mathrm{ep}}(q_{\max})\lambda(q)^\beta\), and squaring both nonnegative sides,
\[
A_\beta\!\left(p^{\mathrm{Ch}}(k,q)\right)
\le
\bigl(K(c)C_{\mathrm{ep}}(q_{\max})\bigr)^2\lambda(q)^{2\beta}
=
C_+(c,q_{\max})
\left(\frac{(1+\sqrt{1-q})^2}{q}\right)^{2\beta}.
\]
% lean: chebyshev_amplification_upper

\item \emph{The two-sided minimax bound.}
The first two displayed inequalities of the statement are exactly the bounds just proved. For the minimax
lower bound, the universal lower estimate obtained from \cref{obj:lem:unbiased-weight-set-nonempty}, \cref{obj:lem:amplification-dual-norm}, and \cref{obj:lem:continuous-chebyshev-endpoint-bound} holds for every \(p\in S_{k,q}\), and the schedule class is
nonempty because \(p^{\mathrm{Ch}}(k,q)\in S_{k,q}\); so the common bound is a lower bound for the nonempty set \(\{A_\beta(p):p\in S_{k,q}\}\), and passing to its infimum gives
\[
M_{\beta,k,q}
\ge
C_-(q_{\max})
\left(\frac{(1+\sqrt{1-q})^2}{q}\right)^{2\beta}.
\]
% lean: minimaxAmplification_lower_of_forall
For the minimax upper bound, that set is bounded below by \(0\) and contains \(A_\beta(p^{\mathrm{Ch}}(k,q))\), so its infimum is at most that value; chaining with the Chebyshev upper bound obtained from \cref{obj:lem:chebyshev-schedule-admissible}, \cref{obj:lem:oversampled-chebyshev-lobatto-norming}, and \cref{obj:lem:continuous-chebyshev-endpoint-bound},
\[
M_{\beta,k,q}
\le
A_\beta\!\left(p^{\mathrm{Ch}}(k,q)\right)
\le
C_+(c,q_{\max})
\left(\frac{(1+\sqrt{1-q})^2}{q}\right)^{2\beta}.
\]
This proves the claimed two-sided minimax estimate.
% lean: minimaxAmplification_le_of_budgeted
\end{enumerate}
\end{proof}
